% This is samplepaper.tex, a sample chapter demonstrating the
% LLNCS macro package for Springer Computer Science proceedings;
% Version 2.21 of 2022/01/12
%
\documentclass[runningheads]{llncs}
%
\usepackage[T1]{fontenc}
% T1 fonts will be used to generate the final print and online PDFs,
% so please use T1 fonts in your manuscript whenever possible.
% Other font encondings may result in incorrect characters.
%
\usepackage{graphicx}
\usepackage{amsmath,amssymb}
% Used for displaying a sample figure. If possible, figure files should
% be included in EPS format.
%
% If you use the hyperref package, please uncomment the following two lines
% to display URLs in blue roman font according to Springer's eBook style:
%\usepackage{color}
%\renewcommand\UrlFont{\color{blue}\rmfamily}
%\urlstyle{rm}
%
\begin{document}
%
\title{Volumetric Gaussian Splatting: Direct Optimization in Voxel Space for Efficient Medical Data Visualization}
%
%\titlerunning{Abbreviated paper title}
% If the paper title is too long for the running head, you can set
% an abbreviated paper title here
%
\author{First Author\inst{1}\orcidID{0000-1111-2222-3333} \and
Second Author\inst{2,3}\orcidID{1111-2222-3333-4444} \and
Third Author\inst{3}\orcidID{2222--3333-4444-5555}}
%
\authorrunning{F. Author et al.}
% First names are abbreviated in the running head.
% If there are more than two authors, 'et al.' is used.
%
\institute{Princeton University, Princeton NJ 08544, USA \and
Springer Heidelberg, Tiergartenstr. 17, 69121 Heidelberg, Germany
\email{lncs@springer.com}\\
\url{http://www.springer.com/gp/computer-science/lncs} \and
ABC Institute, Rupert-Karls-University Heidelberg, Heidelberg, Germany\\
\email{\{abc,lncs\}@uni-heidelberg.de}}
%
\maketitle              % typeset the header of the contribution
%
\begin{abstract}
We present a novel formulation of 3D Gaussian Splatting for volumetric data that departs from the conventional paradigm of multi-view image supervision. Instead of optimizing against rendered RGB images, our approach directly operates in voxel space using target volumes and region-of-interest constraints, enabling a volume-native optimization process. To this end, we introduce a differentiable splat-to-volume rasterization that projects anisotropic Gaussians into a voxel grid and allows direct optimization of spatial and radiometric parameters.

Each Gaussian explicitly encodes a physical intensity value, such as Hounsfield units in computed tomography, enabling a volume-rendering-like representation with post hoc transfer function control. This allows flexible visualization without retraining and establishes a closer connection to classical volume rendering pipelines.

Our method further incorporates volume-driven initialization strategies and stable optimization mechanisms tailored to volumetric supervision. We demonstrate that it enables significantly more efficient rendering, making it particularly suitable for resource-constrained environments such as mixed reality devices.

Overall, our work reframes Gaussian Splatting as a direct volumetric representation and optimization problem, opening new directions for efficient and controllable visualization of medical data.

\keywords{Gaussian Splatting \and Medical Visualization \and Volume Rendering \and Differentiable Rendering \and Volumetric Optimization}
\end{abstract}
%
%
%

\section{Introduction}

Recent advances in medical image analysis have led to highly accurate and increasingly automated segmentation of anatomical structures, driven by methods such as TotalSegmentator \cite{totalsegmentator} and foundation models like SAM \cite{sam}. While these developments enable rich semantic understanding of volumetric data, efficient and high-quality visualization of such data — particularly in immersive environments — remains a significant challenge.

Volume rendering is the de facto standard for visualizing medical data, as it preserves the full volumetric information and allows flexible exploration via transfer functions. However, achieving high visual fidelity typically requires substantial computational resources, which limits its applicability in resource- constrained settings such as mixed reality (MR) devices. Existing solutions often rely on streaming-based approaches, where rendered images or intermediate representations are transmitted from powerful servers to lightweight clients. While effective, these approaches introduce system complexity, latency, and potential failure modes that are undesirable in clinical or interactive scenarios.

At the same time, 3D Gaussian Splatting \cite{kerbl2023gaussiansplatting} has recently emerged as a highly efficient rendering paradigm, enabling real-time visualization of complex scenes with impressive performance characteristics. However, existing formulations are inherently tied to multi-view RGB supervision and are primarily designed for reconstructing scenes from images. This paradigm is fundamentally mismatched with medical imaging, where volumetric ground-truth data is directly available and where accuracy and consistency are critical.

In this work, we propose a reformulation of 3D Gaussian Splatting that directly operates in volumetric space. Instead of relying on image-based supervision, we optimize a set of anisotropic Gaussians against a target volume, guided by region-of-interest constraints. This eliminates the need for intermediate rendering during training and aligns the representation more closely with the underlying data domain.

A key aspect of our approach is an intensity-aware Gaussian representation, where each splat encodes a physical intensity value (e.g., Hounsfield units). This enables a volume-rendering-like behavior in which transfer functions can be modified after training, providing flexibility comparable to classical volume rendering techniques. Furthermore, we introduce a differentiable splat-to-volume rasterization that allows direct optimization in voxel space, as well as volume-driven initialization strategies and stabilization techniques tailored to this setting.

While our method does not aim to fully replace high-quality volume rendering in terms of visual fidelity, it offers a compelling trade-off between quality and efficiency. In particular, it enables fast, lightweight rendering that is well-suited for interactive and immersive applications where computational resources are limited.

In summary, our work bridges the gap between Gaussian Splatting and volumetric medical data by reframing the problem as a direct optimization in voxel space. This opens up new possibilities for efficient, controllable, and domain-aligned visualization of medical volumes.

Our main contributions are as follows:
\begin{itemize}
\item \textbf{Direct volumetric supervision for 3D Gaussian Splatting.} We reformulate 3D Gaussian Splatting by replacing conventional multi-view RGB supervision with direct optimization in voxel space using target volumes and region-of-interest masks.

\item \textbf{Differentiable splat-to-volume rasterization.} We introduce a differentiable rasterization method that projects anisotropic Gaussians into a voxel grid, enabling direct optimization of spatial and radiometric parameters in volumetric space.

\item \textbf{Volume-native Gaussian representation with intensity-aware splats.} Each Gaussian explicitly encodes a physical intensity value (e.g., Hounsfield units), enabling a volume-rendering-like representation where transfer functions can be modified post hoc without retraining.

\item \textbf{Efficient volumetric rendering for resource-constrained devices.} We demonstrate that the proposed representation enables efficient, volume-rendering-like visualization suitable for real-time applications on resource-constrained devices such as mixed reality hardware.
\end{itemize}

\section{Related Work}

\subsection{Gaussian Splatting and Neural Rendering}

Recent advances in neural rendering have introduced highly efficient scene representations, among which 3D Gaussian Splatting has emerged as a particularly impactful approach. The foundational work 3D Gaussian Splatting for Real-Time Radiance Field Rendering represents scenes as sets of anisotropic Gaussians optimized from multi-view image supervision, enabling real-time rendering with high visual fidelity.

Subsequent work has explored multiple extensions of this paradigm. Cinematic Gaussians demonstrates that Gaussian Splatting can achieve near photorealistic visualization of anatomical structures using advanced rendering techniques such as path tracing. Similarly, Relightable 3D Gaussian introduces relightable representations that further reduce the need for traditional volume rendering.

Recent work also investigates more advanced formulations of Gaussian-based rendering. DDGS-CT: Direction-Disentangled Gaussian Splatting for Realistic Volume Rendering and 6DGS: Enhanced Direction-Aware Gaussian Splatting for Volumetric Rendering extend Gaussian Splatting toward volumetric rendering scenarios by incorporating direction-dependent effects and improving fine-detail reconstruction. Other extensions explore alternative primitives such as Triangle Splatting for Real-Time Radiance Field Rendering or hybrid pipelines with explicit clipping support RaRa Clipper. More recent work such as iVR-GS: Inverse Volume Rendering for Explorable Visualization via Editable 3D Gaussian Splatting introduces transfer-function-like behavior by combining multiple Gaussian representations.

Despite these advances, all existing approaches remain fundamentally tied to image-based supervision and view-dependent rendering. Even methods targeting volumetric effects still rely on rendering-based optimization rather than directly leveraging volumetric ground-truth data.



\subsection{Volume Rendering and Splatting Techniques}

Volume rendering is the standard approach for visualizing volumetric data and has been extensively studied in computer graphics. Early work such as Volume Rendering established the foundations for direct volume visualization, while subsequent approaches have focused on improving realism and interactivity. For instance, Exposure Render: An Interactive Photo-Realistic Volume Rendering Framework demonstrates high-quality photorealistic rendering using physically-based techniques.

Splatting-based approaches, such as EWA Volume Splatting, approximate volume rendering using Gaussian kernels and share conceptual similarities with modern Gaussian Splatting methods. More recent discussions, such as Does 3D Gaussian Splatting Need Accurate Volumetric Rendering?, highlight ongoing challenges in achieving both high fidelity and efficiency.

Despite significant progress, classical and modern volume rendering approaches remain computationally demanding, particularly for large scale medical datasets and real time applications.



\subsection{Efficient Medical Visualization and Mixed Reality}


Efficient visualization of volumetric medical data is critical for interactive and immersive applications. Survey works such as Mixed Reality in the Operating Room: A Systematic Review highlight both the potential and limitations of mixed reality in clinical workflows.

Prior work has explored optimized rendering techniques such as A Clinical User Study Investigating the Benefits of Adaptive Volumetric Illumination Sampling, which demonstrate that improved volume rendering can positively impact clinical decision-making. However, even optimized approaches may struggle to meet performance requirements on resource-constrained devices without additional system complexity.


\subsection{Segmentation and Structure-Aware Medical Imaging}


Recent advances in medical image analysis have enabled accurate and large-scale segmentation of anatomical structures. Methods such as nnU-Net, TotalSegmentator, and Segment Anything Model provide detailed structural information directly from volumetric data.

However, most visualization pipelines do not integrate this structural information at the representation level, limiting efficiency and flexibility.


\subsection{Summary}

In summary, Gaussian Splatting has enabled highly efficient rendering but remains fundamentally image-based, even in recent volumetric extensions. At the same time, classical volume rendering provides accurate visualization at high computational cost. Medical imaging, however, offers direct access to volumetric data and structural annotations. Our work bridges this gap by reformulating Gaussian Splatting as a direct volumetric optimization problem, enabling efficient and flexible visualization without intermediate rendering.




\section{Method}

\subsection{Overview}
We reformulate 3D Gaussian Splatting as a direct volumetric optimization problem. Instead of learning a scene representation from calibrated multi-view RGB images, we optimize a set of anisotropic 3D Gaussians directly against a target volume $V$ and a corresponding region-of-interest (ROI) mask $M$. This volume-native formulation removes the need for camera models and image-based supervision, making it particularly suitable for medical and scientific data.

Given a target volume $V$ and an ROI mask $M$, our goal is to learn a compact set of Gaussians that approximates the underlying volumetric signal in relevant regions. Optimization is performed entirely in voxel space using a differentiable splatting formulation.

\subsection{Differences to Standard Gaussian Splatting}
Our method differs fundamentally from standard 3D Gaussian Splatting. The original 3DGS pipeline is designed for novel-view synthesis from calibrated images and optimizes view-dependent color through image-space rendering. In our formulation, there is no camera setup and no image-plane supervision. Instead, the representation is optimized directly against a 3D target volume.

This shift changes every major stage of the pipeline. Initialization is driven by masks rather than by structure-from-motion geometry, rendering is performed into a voxel grid rather than into images, losses are defined in volume space, and densification relies on voxel-space optimization signals instead of visibility heuristics. As a result, Gaussian Splatting becomes a sparse, trainable volumetric proxy rather than purely a rendering primitive for image reproduction.

\subsection{Gaussian Representation}
We represent the volume using a set of anisotropic Gaussians. Each Gaussian is defined by a 3D position, anisotropic scale, rotation, opacity, and a scalar intensity value.

A key design choice of our method is the use of intensity-aware Gaussians. Each primitive explicitly encodes a physical intensity value (e.g., Hounsfield units in CT), allowing the representation to model both structure and appearance in a unified way. This enables transfer functions to be applied or modified after training, similar to classical volume rendering.







\subsection{Initialization}
Initialization is explicitly mask-centric. Instead of starting from sparse-view
geometry or random points, we sample candidate voxels from the ROI mask and
place Gaussian centers inside relevant foreground regions. The initializer
supports both binary masks and soft masks. In the soft-mask case, the method
retains low-confidence but potentially relevant regions by using a permissive
mask threshold, while still focusing sampling on voxels with meaningful mask
support.

To avoid pathological oversampling in dense regions, the initialization stage
uses a coarse per-cell quota before placing points. Sampled voxel coordinates
are then perturbed with sub-voxel jitter and bounded positional noise, after
which invalid samples are corrected if they drift outside the mask support.
This yields continuous initial positions rather than points locked to the
integer voxel lattice.

Initial scales are drawn in voxel units from a configurable range and then
mapped into the normalized volume coordinate system. Opacity is initialized
from mask support, optionally with gamma remapping. If a source intensity
volume is provided, each Gaussian also receives an initial scalar intensity
sampled from the volume. The implementation includes additional boundary-aware
corrections and fallback nearest-voxel sampling to make this process robust
near soft boundaries.

The initializer can further inject structural priors. Initial rotations are
derived from a gradient-based orientation field computed from the intensity
volume. Optionally, Hessian-based structure cues on the mask are used to
estimate vesselness and preferred local directions. These cues can stretch
splats along vessel-like structures, blend structure-derived orientations into
the initial quaternions, and flatten splats near boundaries. This is important
for thin anatomical structures, where isotropic seeds are a poor starting
point.



\subsection{Training and Optimization}

\subsubsection{Splat-to-Volume Formulation}
In order to being able to optimize the gaussian set against the groundtruth volume, the Gaussian set is converted into a dense voxel grid through a differentiable splatting operation. Each Gaussian contributes to nearby voxels based on its spatial extent and orientation, and the final volume is obtained by accumulating all contributions.

\[
\hat{V}(\mathbf{x}) = \sum_{i} \alpha_i , c_i , \mathcal{N}(\mathbf{x} \mid \mu_i, \Sigma_i)
\]

This formulation allows direct comparison between the predicted volume and the target data without intermediate image rendering.

\subsubsection{Training Loop and Loss Functions}

Training is performed directly in voxel space using a dedicated supervision
module that manages the target volume, the mask volume, normalization ranges,
and branch-specific loss computation. The method supports three supervision
modes: density-only mask fitting, intensity-only fitting, and joint
optimization of both branches. In all cases, the losses are evaluated on the
voxel grid rather than on rendered images.

To reduce domination by large empty background regions, supervision is focused
on the mask-defined support. The implementation supports Dice, Tversky, MSE,
and KL losses, which makes the framework suitable for both sparse
segmentation-like targets and continuous scalar reconstruction. In joint mode,
mask and intensity terms are combined with separate weights, and an additional
outside-mask penalty discourages density leakage outside the foreground:
\begin{equation}
\mathcal{L} = \lambda_m \mathcal{L}_m + \lambda_c \mathcal{L}_c + \lambda_o \mathcal{L}_{\mathrm{out}}.
\end{equation}
This compact form is sufficient to capture the training objective without
overcomplicating the presentation.

The training loop includes several mechanisms that are important in practice.
Mixed precision is used for efficiency, and gradient clipping is applied to
avoid instability. Because the point count can exceed the feasible forward-pass
 budget, the implementation optionally renders only an active subset of
Gaussians per iteration while cycling fairly through the full set over time.
Attribute buffers for intensity and opacity can be refreshed periodically from
the source volume or mask, especially for large splats that should remain
aligned with the underlying data.

The topology of the Gaussian set is adapted during training through
densification and pruning. Instead of relying on visibility-based heuristics,
the method accumulates per-point position-gradient norms in voxel space and
periodically clones or splits high-gradient splats while removing low-opacity
ones. This is particularly useful for sparse and filamentary structures, where
fixed topology easily leads to holes or fragmented coverage.



\subsection{Implementation Details}
The implementation is organized around three main components. The training
entry point wires together the Gaussian model, the volume supervisor, the
initialization logic, the optimizer, and export utilities. Initialization is
handled in a dedicated module that performs mask sampling, intensity sampling,
orientation setup, anisotropy injection, and border-aware seeding. Volume
supervision is handled separately and encapsulates target loading, loss
computation, and sampled attribute refresh.

Practical training on full 3D data requires explicit memory management. The
volume loader supports safety downscaling for very large inputs and configurable
storage dtypes such as fp32, fp16, and bf16. The rasterizer supports reduced
internal working grids, batched accumulation, and gradient checkpointing to
trade computation for memory during backpropagation. Together, these choices
make it possible to optimize on larger volumes and higher point counts than a
naive full-resolution implementation would permit.

The code also includes several diagnostics and export features that are
important for research workflows. TensorBoard logging, parameter monitoring,
gradient statistics, and checkpoint saving are integrated into the training
loop. In addition, the Gaussian set can be exported as a sequence of PLY files
throughout training.

The implementation is based on the original 3D Gaussian Splatting code and available at \url{https://github.com/your-repo/3d-gaussian-splatting}.


% ----------------------------

\section{Experiments}

\subsection{Experimental Setup}
We evaluate our method on medical volumetric datasets, focusing on anatomically relevant structures such as vessels and organs. Input data consists of volumetric scans (e.g., CT) together with corresponding segmentation masks.

All volumes are normalized to a common intensity range and processed in voxel space. Gaussian representations are initialized from ROI masks and optimized using the proposed voxel-space formulation.

We compare our method against classical volume rendering approaches and analyze both visual quality and computational efficiency. All experiments are conducted on a standard GPU setup.


\subsection{Baselines}
We compare our approach against the following baselines:

\textbf{Volume Rendering (VR):}
Standard GPU-based volume rendering with transfer functions. This serves as the primary reference for visual fidelity.

\textbf{Optimized Volume Rendering:}
An optimized variant (e.g., adaptive sampling or AVIS-style rendering) representing the current state of efficient medical visualization.

\textbf{3D Gaussian Splatting (Image-Based):}
Standard 3DGS trained via multi-view rendering (if applicable), highlighting the mismatch between image-based and volumetric supervision.


\subsection{Evaluation Metrics}

We evaluate both reconstruction quality and efficiency:

\textbf{Quality Metrics:}
\begin{itemize}
\item Mean Squared Error (MSE) between predicted and target volume
\item Dice score for structure reconstruction
\item Visual comparison of rendered outputs
\end{itemize}

\textbf{Efficiency Metrics:}
\begin{itemize}
\item Rendering time per frame
\item Frames per second (FPS)
\item Memory consumption
\end{itemize}



\subsection{Quantitative Results}

We compare reconstruction accuracy and runtime performance across all methods.

\textbf{Key observation:}
Our method achieves competitive reconstruction quality while significantly reducing rendering cost compared to volume rendering approaches.

\textbf{Table idea:}

%\begin{center}
%\begin{tabular}{lccc}
%Method & MSE ↓ & Dice ↑ & FPS ↑ \
%Volume Rendering & ... & ... & ... \
%Optimized VR & ... & ... & ... \
%Ours & ... & ... & ... \
%\end{tabular}
%\end{center}


\subsection{Qualitative Results}

We visually compare our method to classical volume rendering.

\textbf{Figure idea:}
\begin{itemize}
\item Volume Rendering (reference)
\item Ours (Gaussian Splatting)
\item Difference map (optional)
\end{itemize}

Our method preserves overall anatomical structure while producing a significantly more efficient representation.


\subsection{Efficiency Analysis}

We analyze the trade-off between visual fidelity and rendering performance.

Our approach replaces dense voxel traversal with a sparse set of optimized Gaussians, leading to significantly improved runtime performance, especially for large volumes.

This makes the method particularly suitable for resource-constrained environments such as mixed reality devices.


\subsection{Ablation Study}

We analyze the impact of key components of our method:

\textbf{Without ROI supervision:}
Leads to degraded performance due to background dominance.

\textbf{Without intensity learning:}
Prevents transfer-function-like behavior and reduces visual flexibility.

\textbf{Without densification/pruning:}
Results in inefficient or incomplete representations.

\textbf{Without anisotropic Gaussians:}
Reduces reconstruction quality for thin structures.


\subsection{Discussion}

Our results demonstrate that direct volumetric optimization enables a new trade-off between quality and efficiency.

While classical volume rendering still achieves higher visual fidelity in some cases, our method provides a significantly more efficient alternative that remains suitable for interactive and real-time applications.

This makes it particularly promising for medical visualization in mixed reality environments, where performance constraints are critical.


\section{Conclusion}
We presented a volume-supervised formulation of Gaussian Splatting that replaces multi-view image supervision with direct optimization in voxel space. By fitting anisotropic 3D Gaussians directly to a target volume and ROI mask, the method turns Gaussian Splatting into a volume-native representation that is better aligned with medical and scientific data.

The proposed pipeline combines mask-driven and structure-aware initialization, differentiable splat-to-volume rendering, flexible density and intensity supervision, and a set of practical training mechanisms for memory efficiency and stability. Beyond optimization itself, the learned Gaussian set can be exported as an interpretable proxy for downstream visualization and analysis.

Overall, this work establishes a concrete foundation for bridging efficient Gaussian-based rendering with direct volumetric data supervision. We believe this perspective opens a promising line of research toward sparse, controllable, and application-aware 3D representations for medical visualization.

\begin{credits}
\subsubsection{\ackname} A bold run-in heading in small font size at the end of the paper is
used for general acknowledgments, for example: This study was funded
by X (grant number Y).

\subsubsection{\discintname}
It is now necessary to declare any competing interests or to specifically
state that the authors have no competing interests. Please place the
statement with a bold run-in heading in small font size beneath the
(optional) acknowledgments\footnote{If EquinOCS, our proceedings submission
system, is used, then the disclaimer can be provided directly in the system.},
for example: The authors have no competing interests to declare that are
relevant to the content of this article. Or: Author A has received research
grants from Company W. Author B has received a speaker honorarium from
Company X and owns stock in Company Y. Author C is a member of committee Z.
\end{credits}



\bibliographystyle{splncs04}
\bibliography{bibl}

\end{document}
